% =============================================================================
%  00_front.tex —— 书名页、译序、体例说明
% =============================================================================

\begin{titlepage}
\centering
\vspace*{2.6cm}

{\Huge\bfseries 给数学家的导出范畴\par}

\vspace{1.1cm}
{\Large R.\,P.~Thomas\quad 著\par}
\vspace{0.35cm}
{\large 中文译注本\par}

\vspace{2.2cm}
{\large 底本\par}
\vspace{0.3cm}
{\large\itshape Derived categories for the working mathematician\par}
\vspace{0.5cm}
{\large arXiv:math/0001045\quad（2000 年 1 月 7 日；2001 年 10 月 17 日修订）\par}
\vspace{0.3cm}
{\large 会议版：\emph{Winter School on Mirror Symmetry, Vector Bundles\\[2pt]
and Lagrangian Submanifolds}（Harvard, 1999）\par}

\vfill
{\large 数译（Shuyi）\quad 译注\par}
\vspace{0.3cm}
{\large 二〇二六年十一月\par}
\vspace{1.2cm}
{\footnotesize
底本权利归原作者与出版方。本译本为学习与研究用途的演绎作品，\\
不替代原文；如权利人提出要求，可随时撤下。\par}
\end{titlepage}

\thispagestyle{plain}
\section*{译序}
\addcontentsline{toc}{section}{译序}

这篇讲稿有一个著名的标题。Thomas 借 MacLane 的《Categories for the working
mathematician》拟出“给数学家的导出范畴”，开篇第一句却是为标题道歉：讲稿真正的
听众，至少有一半是物理学家。

那是一九九九年一月，哈佛的镜像对称冬季学校。Kontsevich 的同调镜像对称纲领把
三角范畴推到了弦论与几何的主干道上，D-膜的数学描述又让它变成日常语言。可对刚
进门的人，导出范畴的名声并不好：定义绕、文献旧、关键文章用法文写。Thomas 的
处理办法不是把定义抄一遍，而是先回答一个问题——我们为什么需要它？

他的答案极简：\emph{把层换成复形，把拟同构当成同构}。全文十三页，没有一条
需要查书的证明，却把 Ext、Tor、超上同调、Leray 谱序列、Hodge 分解统统还原成
同一个动作。读完再回头看 Hartshorne 的《Residues and Duality》，那些吓人的
双复形论证会突然显得只是技术细节。

正因为它是“最短路径”，译本的价值不在于逐字对应，而在于把跳过的台阶补出来。
本译本做了四件事。

\begin{enumerate}[leftmargin=1.6em,itemsep=0.2em]
\item \textbf{全译}。正文六节、摘要、参考文献全部译出；公式编号、节号与
      arXiv 版一致，方便对照原文逐式核查。
\item \textbf{译注}。凡原文一句话带过、而读者可能卡住的地方（例如“双重对偶
      不是原来的同调”、“谱序列的指标怎么排”、“为什么锥能同时充当核与余核”），
      加橙色译注框，并在末句标明依据。
\item \textbf{概念卡片}。拟同构、局部化、三角范畴、导出函子等核心概念，
      用紫色卡片给出可供查证的标准表述，与正文互相印证。
\item \textbf{校勘}。原文有若干处笔误与表述不确（节号错位、一个同调不同的
      例子、Dolbeault 复形与全纯 de Rham 复形混称、谱序列记号残缺等），
      逐一在正文以红色校勘框标出，并在附录 B 给出判定、依据与建议改法。
      正文不擅自改动原文的数学内容。
\end{enumerate}

\subsection*{三条读法}

\begin{itemize}[leftmargin=1.6em,itemsep=0.2em]
\item \textbf{拓扑路线}（最快）：§\ref{sec:cx} 的单纯复形与 Whitehead 定理
      $\to$ §\ref{sec:cone} 的锥与 Dold--Puppe 构造 $\to$ §\ref{sec:cone}
      末尾的三角范畴。不需要层论。
\item \textbf{代数几何路线}：§\ref{sec:cx} 的除子与 Koszul 复形
      $\to$ §\ref{sec:qi} 的拟同构 $\to$ §\ref{sec:df} 的导出函子与
      谱序列。读这部分时手边放一个光滑曲面上的除子例子最省力。
\item \textbf{范畴论路线}：§\ref{sec:cat} 的加性/阿贝尔范畴与 $D^b(\A)$
      的局部化构造 $\to$ §\ref{sec:cone} $\to$ §\ref{sec:df}。
      这条路最抽象，但也是唯一能说明“为什么不能取核与余核”的路。
\end{itemize}

附录 D 是译者写的精读指南：一张论证地图、每节的关键跳跃、以及若干自测题。
它不替代正文，只在读第二遍时使用。

\subsection*{体例}

\begin{itemize}[leftmargin=1.6em,itemsep=0.2em]
\item 术语首次出现时给中英对照，之后只用中文；术语总表见附录 A。
\item 定理、定义、例、注沿用原文的共享计数器与按节编号（如“定义 4.1”），
      与 arXiv 版一一对应。
\item 译注框、概念卡片、校勘框均可通过导言区的开关一键隐藏
      （\verb|\showyizhufalse| 等），隐藏后文字以内文排印。
\item 原文的强调（原稿用无衬线粗体）在译本中作\textbf{粗体}。
\end{itemize}

\subsection*{致谢}

感谢原作者 R.\,P.~Thomas 写下这篇讲稿，并把它公开在 arXiv 上。译文中凡属
译者的补充与判断，均已标明；如有误读，责任在译者。

\vspace{0.5em}
\begin{flushright}
数译\quad 二〇二六年十一月
\end{flushright}

\newpage
